
For this it is clearly necessary that @@M0@@ be catenary, and sufficient that @@M1@@ be biequidimensional (EGA @@M2@@ 14.3.3). Nevertheless, there exist noetherian schemes of finite dimension, catenary and equidimensional, having no such functions @@M3@@. To see this, one can modify the example (EGA IV 5.6.11) as follows: take two copies @@M4@@ and @@M5@@ of the spectrum of the ring @@M6@@ of loc. cit., and denote by @@M7@@ and @@M8@@ (@@M9@@) the closed points of @@M10@@ (@@M11@@) corresponding to the maximal ideals @@M12@@ and @@M13@@ of @@M14@@, of heights @@M15@@ and @@M16@@ respectively. Recall that the residue fields of @@M17@@ and @@M18@@ are isomorphic. Thus one can construct a scheme @@M21@@ by gluing @@M19@@ and @@M20@@, identifying @@M22@@ with @@M23@@ and @@M24@@ with @@M25@@. One can then identify @@M26@@ and @@M27@@ with the two irreducible components of @@M28@@, and then, inside @@M29@@, @@M30@@ has dimension zero. Thus @@M31@@ is a noetherian, catenary and equidimensional scheme (of dimension @@M32@@). If (ii bis) held on @@M33@@, one would have
@@M34@@
and, on the other hand,
@@M35@@
giving a contradiction.

